\cotefile{68}
\scanpage{33}
% archive foliation: 32
\[
\left\{\begin{aligned}
G\cap(C\setminus C_0)&=\unsure{G^\times|G^*}\\
P\cap(C\setminus C_0)&=\unsure{P^*|P^\times}\\
C\setminus(G\cup P)&=N
\end{aligned}\right.
\qquad\text{(positions nulles, i.e. gagnantes pour aucun des deux joueurs)}
\]
Une position $x$ sera encore dite \uline{gagnante} (resp. \uline{perdante}) si $x\in G$ (resp. $x\in P$) --- cette terminologie coïncide avec la précédente si $x\notin C_0$. On aura, posant
\[
\deleted{\unsure{G(j)=\{x\in}}
\]
\[
\hat C(j)=\{x\in C\mid\unsure{\hat\alpha}(x)=j\}
\]
(de sorte que $C$ \unsure{apparait|apparaît} comme réunion disjointe des deux ens. $\hat C(j)$, $j\in J$),
\[
\left\{\begin{aligned}
G(j)\cap\hat C(j)&=\deleted{\unsure{\hat\alpha}}\ \hat C(j)\cap G\\
G(j)\cap\hat C(\bar j)&=\hat C(\bar j)\cap P
\end{aligned}\right.
\]
donc \inserted{les} $G(j)$ $(j\in J)$ se \unsure{reconstituent} à partir de $\unsure{\hat\alpha}$ et des ensembles $G,P$.

On aura, pour $x\in C\setminus C_0$
\[
\begin{aligned}
x\in G&\Longleftrightarrow\exists y\in R(x)\ \inserted{\text{tel que}}\ y\in P
&&\text{i.e. }R(x)\cap P\ne\varnothing\\
x\in P&\Longleftrightarrow\forall y\in R(x),\ \text{on a }y\in G
&&\text{i.e. }R(x)\unsure{\in|\subset}G
\end{aligned}
\]
i.e.
\[
\begin{aligned}
G\cap(C\setminus C_0)\deleted{\unsure{G}}&=\{x\in C\setminus C_0\mid R(x)\cap P\ne\varnothing\}\\
P\cap(C\setminus C_0)&=\{x\in C\setminus C_0\mid R(x)\subset G\}
\end{aligned}
\]
